% =============================================================================
%  DIAGRAMA DE ARQUITECTURA — Figura 4.1
%
%  PREÁMBULO: añade estas líneas a main.tex, junto a los demás \usepackage
%
%    \usepackage{tikz}
%    \usetikzlibrary{positioning, arrows.meta, fit, backgrounds, calc}
%
%  USO: sustituye el bloque \framebox de la sección 4.1 por todo lo que sigue
%
%  NOTA: el código evita deliberadamente los caracteres < y >, que babel con
%  la opción 'spanish' convierte en caracteres activos (atajo de «»), lo que
%  rompería la sintaxis habitual de TikZ (-> y >=).
% =============================================================================

\begin{figure}[H]
\centering
\begin{tikzpicture}[
  font=\small,
  % ── Estilos ──────────────────────────────────────────────────────────────
  fuente/.style={
    draw=azultfm!60, fill=azultfm!8, rounded corners=2pt,
    minimum width=2.5cm, minimum height=0.85cm, align=center,
    inner sep=3pt, font=\footnotesize},
  proceso/.style={
    draw=gray!70, fill=white, rounded corners=2pt,
    minimum width=3.2cm, minimum height=0.85cm, align=center,
    inner sep=3pt, font=\footnotesize},
  almacen/.style={
    draw=gray!70, fill=gray!12, rounded corners=2pt,
    minimum width=2.6cm, minimum height=0.85cm, align=center,
    inner sep=3pt, font=\footnotesize},
  compartido/.style={
    draw=orange!80!black, fill=orange!12, rounded corners=2pt,
    line width=0.9pt,
    minimum width=3.4cm, minimum height=1.0cm, align=center,
    inner sep=3pt, font=\footnotesize},
  salida/.style={
    draw=green!45!black, fill=green!10, rounded corners=2pt,
    minimum width=3.2cm, minimum height=0.85cm, align=center,
    inner sep=3pt, font=\footnotesize},
  flecha/.style={-{Stealth[length=2.2mm]}, draw=gray!65, line width=0.7pt},
  flechacomp/.style={-{Stealth[length=2.2mm]}, draw=orange!80!black,
    line width=0.9pt, dashed},
  etiqueta/.style={font=\scriptsize\itshape, text=gray!60!black},
  capa/.style={rounded corners=3pt, draw=gray!35, line width=0.6pt},
  titulocapa/.style={font=\footnotesize\bfseries, text=gray!55!black}
]

% ═══════════════════════════════════════════════════════════════════════════
% CAPA DE DATOS
% ═══════════════════════════════════════════════════════════════════════════
\node[fuente] (ree)     {REE\\\tiny REData};
\node[fuente, right=0.45cm of ree]    (aemet)  {AEMET\\\tiny OpenData};
\node[fuente, right=0.45cm of aemet]  (cal)    {Calendario\\\tiny\texttt{holidays}};

\node[proceso, below=0.9cm of aemet] (ingesta) {Ingesta y limpieza\\\tiny imputación escalonada};
\node[almacen, below=0.75cm of ingesta] (parquet) {\texttt{dataset.parquet}\\\tiny 26.304 registros};

\draw[flecha] (ree)   -- ([xshift=-1.1cm]ingesta.north);
\draw[flecha] (aemet) -- (ingesta.north);
\draw[flecha] (cal)   -- ([xshift=1.1cm]ingesta.north);
\draw[flecha] (ingesta) -- (parquet);

% ═══════════════════════════════════════════════════════════════════════════
% CAPA DE MODELOS
% ═══════════════════════════════════════════════════════════════════════════
\node[compartido, below=1.1cm of parquet] (feat)
  {\texttt{prediccion.py}\\\tiny construcción de variables};

\node[proceso, below left=0.95cm and -0.6cm of feat] (entrena)
  {Entrenamiento\\\tiny XGBoost};
\node[proceso, below right=0.95cm and -0.6cm of feat] (valida)
  {Validación\\\tiny origen rodante};

\node[almacen, below=0.95cm of entrena] (joblib)
  {\texttt{modelo.joblib}};
\node[almacen, below=0.95cm of valida] (metricas)
  {Métricas\\\tiny MAPE 1,34\,\%};

\draw[flecha] (parquet) -- (feat);
\draw[flecha] (feat.south) -- (entrena.north);
\draw[flecha] (feat.south) -- (valida.north);
\draw[flecha] (entrena) -- (joblib);
\draw[flecha] (valida) -- (metricas);


% ═══════════════════════════════════════════════════════════════════════════
% CAPA DE APLICACIÓN
% ═══════════════════════════════════════════════════════════════════════════
\node[proceso, right=4.2cm of feat] (vivo)
  {Datos en tiempo real\\\tiny bloques de 5 días};
\node[proceso, below=0.95cm of vivo] (predice)
  {Predicción 24--48 h\\\tiny esquema recursivo};
\node[salida, below=0.95cm of predice] (ui)
  {Aplicación Streamlit\\\tiny franjas y recomendación};

\draw[flecha] (vivo) -- (predice);
\draw[flecha] (predice) -- (ui);
% El modelo alimenta la predicción: se encamina por debajo para no
% cruzar los nodos de la capa de modelos
\draw[flecha] (joblib.south) -- ++(0,-0.45) coordinate (ruta)
  -| ([xshift=-0.8cm]predice.west) -- (predice.west);

% ── Código compartido: la relación que da sentido al diseño ───────────────
\draw[flechacomp] (feat.east) -- node[etiqueta, above, pos=0.5]
  {mismo código} (vivo.west);

% ── Fuentes externas que alimentan la capa de aplicación ──────────────────
\draw[flecha, dotted] (cal.north) .. controls +(0,0.9) and +(0,2.6) ..
  node[etiqueta, pos=0.55, right, xshift=2pt] {tiempo real} (vivo.north);

% ═══════════════════════════════════════════════════════════════════════════
% RECUADROS DE CAPA
% ═══════════════════════════════════════════════════════════════════════════
\begin{scope}[on background layer]
  \node[capa, fill=azultfm!4, fit=(ree)(cal)(parquet),
        inner sep=0.30cm] (caja1) {};
  \node[capa, fill=orange!4, fit=(feat)(entrena)(valida)(joblib)(metricas),
        inner sep=0.30cm] (caja2) {};
  \node[capa, fill=green!4, fit=(vivo)(ui),
        inner sep=0.30cm] (caja3) {};
\end{scope}

\node[titulocapa, anchor=south west] at ([yshift=1.5pt]caja1.north west)
  {CAPA DE DATOS};
\node[titulocapa, anchor=south west] at ([yshift=1.5pt]caja2.north west)
  {CAPA DE MODELOS};
\node[titulocapa, anchor=south west] at ([yshift=1.5pt]caja3.north west)
  {CAPA DE APLICACIÓN};

\end{tikzpicture}

\caption[Arquitectura general del sistema]{Arquitectura general del sistema
propuesto. El módulo \texttt{prediccion.py}, destacado en naranja, es compartido
por el procedimiento de entrenamiento y por la aplicación, lo que garantiza que
las variables se construyan de forma idéntica en ambos contextos}
\label{fig:arquitectura}
\end{figure}
